% !TEX root = main.tex
\section{Back Trace the Collatz Sequence}\label{sec:back-trace-the-collatz-sequence}
\cref{theorem:any-presequence} tell us we can also trace back the Collatz sequence in $\mathbb{Q}$.
But we can only go back under a finite reduces given.
For an infinite tracing back, we need to redefine the summation to give it directionality.
\begin{definition}
    Let's modified the definition of $\sum$ operator. The sum $S$ of $s$ from $m$ to $n$ is defined as
    \[ S = \left\{\begin{aligned}
    &\sum_{i=m}^{n-1}s_i, &&n > m, \\
    &0, &&n = m, \\
    -&\sum_{i=n}^{m-1}s_i, &&n < m.
\end{aligned}\right. \]
\end{definition}
With this we have
\begin{proposition}
    $\mathcal{A}_{n}^{m} = -\mathcal{A}_{m}^{n}$.
\end{proposition}
\begin{proposition}
    $\alpha^{n-m} \cdot \mathfrak{A}_{m}^{n} + \beta^{\mathcal{A}_m^n} \cdot \mathfrak{A}_{n}^{m} = 0$ and $\mathfrak{a}_{m}^{n} + \beta^{\mathcal{A}_m^n} \cdot \mathfrak{a}_{n}^{m} = 0$.
\end{proposition}
It is easy to verify that those propositions of Collatz accumulation and complement still holds.
And with these definitions, we have
\begin{theorem}~\label{theorem:final-inductive-formula}
Given $a \in \widehat{\mathbb{Q}}_{{\beta}}$ and $\{v_{i}\}_{i=-\infty}^{\infty} \subset \mathbb{N}$ that extends from the Collatz reduces of $a$ and $\forall i, j \in \mathbb{Z}$, there is $\beta \nmid \mathfrak{A}_{i}^{j}$.
Let $a_0 = a$, then $\forall m, n \in \mathbb{Z}$, there still is
\[ a_n \cdot \beta^{\mathcal{A}_{n}^{m}} = a_m \cdot\alpha^{n-m} + \gamma\cdot\mathfrak{A}_{n}^{m}. \]
\end{theorem}

\begin{corollary}
    $a_n \cdot \beta^{\mathcal{A}_{n}^{m}} = \alpha^{n-m} \cdot\left(a_m + \frac{\gamma}{\alpha}\cdot\alpha^{m}\cdot\mathfrak{a}_{n}^{m}\right)$.
\end{corollary}
\begin{corollary}
    If $\forall i$ that $k \ge i > 0$, there is $v_{-i} = t$, denote $\sigma = \frac{\gamma}{\alpha-\beta^{t}}$, then we have
    $(a_{-k} + \sigma) \cdot \alpha^k  = \beta^{kt} \cdot (a_0 + \sigma)$.
\end{corollary}
Equation $(a_{-k} + \sigma) \cdot \alpha^k  = \beta^{kt} \cdot (a_0 + \sigma)$ has only one form of solution:
\[\left\{\begin{aligned}
    a_{0} &= t \cdot \alpha^k - \sigma, \\
    a_{-k} &= t \cdot \beta^{k\Delta} - \sigma,
\end{aligned}\right.\]
where $t \in \mathbb{Q}$.
To ensure both $a_{-k}$ and $a_0$ are integers, we require
\[ t \in \{s\cdot\sigma+n \mid n \in \mathbb{N}, s \equiv \alpha^{-k} \equiv \beta^{-kt} \pmod{\alpha - \beta^{t}}\}. \]
Unfortunately, in this case, not all solutions of the equation guarantee that $\beta \nmid a_0$.
